% GENERATED FILE. Edit canonical lessons, then run npm run generate:books.
% canonical-source-hash: fnv1a64:b2aa38ccd0b4cd9c
% canonical-lessons: HI-C210-paratha, HI-C210-khicri, HI-C210-halva, HI-C210-laddu, HI-C210-achar

\chapter{A paratha, Khichri, Halwa, A laddu, Pickle}
\label{ch:hi-a-paratha-khichri-halwa-a-laddu-pickle}
\hlchaptermodality{210}

\begin{chapteropening}
\textbf{By the end of this chapter:} \emph{I can say a paratha, khichri, halwa, a laddu and pickle.}

Complete the last lesson of chapter 210: i can say a paratha, khichri, halwa, a laddu and pickle.
\end{chapteropening}

\section[\texorpdfstring{\hi{पराठा} (parāṭhā)}{पराठा (parāṭhā)}]{\hi{पराठा} (parāṭhā) — a paratha}

\label{lesson:HI-C210-paratha}



\begin{quote}
Before the new one: say the Hindi for a lamp, a light, then the Hindi for bread (flatbread).
\end{quote}

\subsection*{You'll want to know: \hi{पराठा}}
\textbf{\hi{पराठा}} — \emph{parāṭhā} — \textquotedblleft{}a paratha\textquotedblright{}.

\begin{grammarlens}[title={What you've built}]
One more word for this chapter.
\end{grammarlens}

\subsection*{Your turn}
\begin{itemize}
\raggedright
  \item \emph{Say it:} \emph{parāṭhā}
  \item \emph{Say it:} \emph{parāṭhā}, once more
  \item \emph{Say it:} say \emph{parāṭhā}
  \item \emph{Recall:} say the Hindi for a lamp, a light, then the Hindi for bread (flatbread), then say \emph{parāṭhā} again
\end{itemize}

\begin{tcolorbox}[breakable,colback=teal!4,colframe=teal!35!black,title={Before you move on}]
What does \hi{पराठा} mean? (\textquotedblleft{}A paratha\textquotedblright{}.) Say it once more.
\end{tcolorbox}
\section[\texorpdfstring{\hi{खिचड़ी} (khicṛī)}{खिचड़ी (khicṛī)}]{\hi{खिचड़ी} (khicṛī) — khichri (rice and lentils)}

\label{lesson:HI-C210-khicri}



\begin{quote}
Before the new one: say the Hindi for bread (flatbread), then the Hindi for a paratha.
\end{quote}

\subsection*{You'll want to know: \hi{खिचड़ी}}
\textbf{\hi{खिचड़ी}} — \emph{khicṛī} — \textquotedblleft{}khichri (rice and lentils)\textquotedblright{}.

\begin{grammarlens}[title={What you've built}]
One more word for this chapter.
\end{grammarlens}

\subsection*{Your turn}
\begin{itemize}
\raggedright
  \item \emph{Say it:} \emph{khicṛī}
  \item \emph{Say it:} \emph{khicṛī}, once more
  \item \emph{Say it:} say \emph{khicṛī}
  \item \emph{Recall:} say the Hindi for bread (flatbread), then the Hindi for a paratha, then say \emph{khicṛī} again
\end{itemize}

\begin{tcolorbox}[breakable,colback=teal!4,colframe=teal!35!black,title={Before you move on}]
What does \hi{खिचड़ी} mean? (\textquotedblleft{}Khichri (rice and lentils)\textquotedblright{}.) Say it once more.
\end{tcolorbox}
\section[\texorpdfstring{\hi{हलवा} (halvā)}{हलवा (halvā)}]{\hi{हलवा} (halvā) — halwa (a sweet)}

\label{lesson:HI-C210-halva}



\begin{quote}
Before the new one: say the Hindi for a paratha, then the Hindi for khichri.
\end{quote}

\subsection*{You'll want to know: \hi{हलवा}}
\textbf{\hi{हलवा}} — \emph{halvā} — \textquotedblleft{}halwa (a sweet)\textquotedblright{}.

\begin{grammarlens}[title={What you've built}]
One more word for this chapter.
\end{grammarlens}

\subsection*{Your turn}
\begin{itemize}
\raggedright
  \item \emph{Say it:} \emph{halvā}
  \item \emph{Say it:} \emph{halvā}, once more
  \item \emph{Say it:} say \emph{halvā}
  \item \emph{Recall:} say the Hindi for a paratha, then the Hindi for khichri, then say \emph{halvā} again
\end{itemize}

\begin{tcolorbox}[breakable,colback=teal!4,colframe=teal!35!black,title={Before you move on}]
What does \hi{हलवा} mean? (\textquotedblleft{}Halwa (a sweet)\textquotedblright{}.) Say it once more.
\end{tcolorbox}
\section[\texorpdfstring{\hi{लड्डू} (laḍḍū)}{लड्डू (laḍḍū)}]{\hi{लड्डू} (laḍḍū) — a laddu (a sweet)}

\label{lesson:HI-C210-laddu}



\begin{quote}
Before the new one: say the Hindi for khichri, then the Hindi for halwa.
\end{quote}

\subsection*{You'll want to know: \hi{लड्डू}}
\textbf{\hi{लड्डू}} — \emph{laḍḍū} — \textquotedblleft{}a laddu (a sweet)\textquotedblright{}.

\begin{grammarlens}[title={What you've built}]
One more word for this chapter.
\end{grammarlens}

\subsection*{Your turn}
\begin{itemize}
\raggedright
  \item \emph{Say it:} \emph{laḍḍū}
  \item \emph{Say it:} \emph{laḍḍū}, once more
  \item \emph{Say it:} say \emph{laḍḍū}
  \item \emph{Recall:} say the Hindi for khichri, then the Hindi for halwa, then say \emph{laḍḍū} again
\end{itemize}

\begin{tcolorbox}[breakable,colback=teal!4,colframe=teal!35!black,title={Before you move on}]
What does \hi{लड्डू} mean? (\textquotedblleft{}A laddu (a sweet)\textquotedblright{}.) Say it once more.
\end{tcolorbox}
\section[\texorpdfstring{\hi{अचार} (achār)}{अचार (achār)}]{\hi{अचार} (achār) — pickle}

\label{lesson:HI-C210-achar}



\begin{quote}
Before the new one: say the Hindi for halwa, then the Hindi for a laddu.
\end{quote}

\subsection*{You'll want to know: \hi{अचार}}
\textbf{\hi{अचार}} — \emph{achār} — \textquotedblleft{}pickle\textquotedblright{}.

\begin{grammarlens}[title={What you've built}]
One more word for this chapter.
\end{grammarlens}

\subsection*{Your turn}
\begin{itemize}
\raggedright
  \item \emph{Say it:} \emph{achār}
  \item \emph{Say it:} \emph{achār}, once more
  \item \emph{Say it:} say \emph{achār}
  \item \emph{Recall:} say the Hindi for halwa, then the Hindi for a laddu, then say \emph{achār} again
\end{itemize}

\begin{tcolorbox}[breakable,colback=teal!4,colframe=teal!35!black,title={Before you move on}]
What does \hi{अचार} mean? (\textquotedblleft{}Pickle\textquotedblright{}.) Say it once more.
\end{tcolorbox}
